\documentclass{beamer}
\usepackage[utf8]{inputenc} 
\title{Elvana Gjata}
\author{Kamil Schilling}
\institute{UWM}
\date{\today}
\usepackage{amsfonts}
\usepackage[MeX]{polski}
\begin{document}
\frame{\titlepage}



\begin{frame}
\frametitle{Spis Treści}
\tableofcontents
\end{frame}

\section{Życiorys Piosenkarki}
\begin{frame}{Życiorys}
\begin {itemize}
\item Pochodzi z rodziny artystycznej (jej starsza siostra Migena jest piosenkarką). Po ukończeniu szkoły średniej rozpoczęła studia w Akademii Sztuk w Tiranie.
Zadebiutowała na scenie jako dziecko, śpiewając w stołecznym Pałacu Pionierów. W 2001 na Festiwalu Młodych Głosów zdobyła trzecią nagrodę (wykonała utwór „Të këndojmë së bashku”). Kolejnym jej sukcesem był występ w programie telewizyjnym Szukamy gwiazdy, w którym doszła do finału wykonując piosenkę z repertuaru Laury Pausini. Od tego momentu zaczęła regularnie pojawiać się w programach telewizji albańskiej.
\end {itemize}
\end{frame}


\begin{frame}{Zdjęcie}

\begin{figure}[here]
\begin{center}
\includegraphics[scale=0.4]{elvana.jpg}
\end{center}
\end{figure}

\end{frame}




\begin{frame}{Życiorys cd.}
\begin {itemize}
\item<1> W 2003 wzięła udział w najbardziej prestiżowym albańskim konkursie dla piosenkarzy, organizowanym pod patronatem Radia i Telewizji. Zaśpiewała w konkursie utwór „Prane teje”, z własnym tekstem, który znalazł się wśród kompozycji wyróżnionych na festiwalu.
\pause
\item<1-3> Współpraca z producentem Mumajesim doprowadziła do wydania w 2007 debiutanckiego albumu Mamës, a także wideoklipu promującego płytę. Rok później ukazał się kolejny album piosenkarki Vetës. Współpracuje z muzykiem Flori Mumajesi.
\item<3> W 2012 wystąpiła w trzeciej edycji programu Dancing with the Stars w parze z Gerdim, odnosząc w nim zwycięstwo.
\end {itemize}
\end {frame}

\section{Dyskografia}
\begin{frame}{Dyskografia}
\begin {itemize}

\item Albumy:

-- 2007: Mamës

-- 2008: Vetës

-- 2011: Afër Larg
\pause

\item Single

--2005 : Te Kam Xhan

--2007 : Pranë Teje

--2007 : Zjarri Që Ti Ndez

--2008 : Ku Jeton Dashuria

--2008 : Ti Ke Ikur

--2009 : Disiplinë

--2009 : Nuk Janë Më



\end {itemize}
\end {frame}

\section{Tekst piosenki}
\begin{frame}{Piosenka}

Tytuł: Me Ty
\newline

Një milion arsye 

Po të kisha unë 

Nuk do të doje kurr 

Si doja dikur . 

S'munde të më thoje fjalën 

Lamtumirë 

Se ti ike shpejt 

për një tjetër jetë ! 

Uuuuuuuhh! 

Ti nuk je aty se .. 

Nuk je si ti e do 

\end {frame}

\begin{frame}{Piosenka cd.}
E une zgjohem në mëngjes 

Sikur jam me ty 

Dërgoj lutje lart në qiell 

Por nuk vin nga ty 

Dhe sa herë kujtimet gjej 

Veten e gënjej 

Unë e gënjej 

Por ty s'të kthej ! 

\end {frame}
\begin{frame}{Piosenka cd.}
Reff : (2x) 

Dhe une zgjohem në mëngjes 

Sikur jam me ty 

Dërgoj lutje lart në qiell 

Por nuk vin nga ty 

Dhe sa herë kujtimet gjej 

Veten e gënjej 

Unë e gënjej 

Por ty s'të kthej ! 

\end {frame}

\section{Koniec}
\begin{frame}{Koniec informacji}
\begin {itemize}
\item To koniec
\end {itemize}
\end {frame}

\section{Bibliografia}
\begin{frame}{Bibliografia}
Bibliografia:
\begin{thebibliography}{9}
\bibitem{deni}
--Wikipedia, Elvana Gjata
\end{thebibliography}
\end{frame}

\end{document}